\usepackage[a4paper,top=20mm,bottom=22mm,left=19mm,right=19mm,headheight=25pt]{geometry}
\usepackage[table]{xcolor}
\usepackage{amsmath,amssymb,mathtools,bm}
\usepackage{siunitx}
\usepackage{booktabs,tabularx,array,multicol,longtable}
\usepackage{enumitem}
\usepackage{graphicx}
\usepackage{float}
\usepackage{caption}
\usepackage{tikz}
\usepackage{pgfplots}
\pgfplotsset{compat=1.18}
\usetikzlibrary{arrows.meta,positioning,calc,angles,quotes,patterns,patterns.meta,
  decorations.pathmorphing,decorations.markings,intersections,shapes.geometric,
  backgrounds,3d,fit,spy}
\usepackage[most]{tcolorbox}
\usepackage{fancyhdr}
\usepackage{hyperref}
\usepackage{bookmark}
\usepackage{xparse}

\definecolor{Ink}{HTML}{202326}
\definecolor{DeepBlue}{HTML}{294E69}
\definecolor{Blue}{HTML}{47687D}
\definecolor{PaleBlue}{HTML}{F3F5F6}
\definecolor{Gold}{HTML}{6E777D}
\definecolor{PaleGold}{HTML}{F6F6F5}
\definecolor{RuleGray}{HTML}{D3D7DA}
\definecolor{TextGray}{HTML}{61686D}
\definecolor{SoftGray}{HTML}{F6F6F6}
\color{Ink}

\sisetup{per-mode=symbol,inter-unit-product=\ensuremath{{}\cdot{}},exponent-product=\times}
\hypersetup{
  colorlinks=true,linkcolor=DeepBlue,urlcolor=Blue,
  bookmarksdepth=2,
  pdftitle={高中物理竞赛一轮讲义},
  pdfsubject={数学工具、力学、热学、电磁学、光学与近代物理的一轮竞赛讲义},
  pdfauthor={物理竞赛一轮资料编写组}
}

\setlength{\parindent}{2em}
\setlength{\parskip}{0.18em}
\setlist{nosep,leftmargin=2em}
\setcounter{secnumdepth}{2}
\setcounter{tocdepth}{1}
\allowdisplaybreaks[2]
\emergencystretch=2em

\ctexset{
  chapter={format=\raggedright\LARGE\bfseries\color{Ink},beforeskip=0pt,afterskip=18pt},
  section={format=\Large\bfseries\color{DeepBlue},beforeskip=1.8ex,afterskip=1.0ex},
  subsection={format=\large\bfseries\color{Ink},beforeskip=1.5ex,afterskip=0.8ex},
  subsubsection={format=\normalsize\bfseries\color{Ink}}
}

\pagestyle{fancy}
\fancyhf{}
\fancyhead[L]{\small\color{TextGray}高中物理竞赛一轮讲义}
\fancyhead[R]{\small\color{TextGray}\nouppercase{\leftmark}}
\fancyfoot[C]{\small\color{TextGray}\thepage}
\renewcommand{\headrulewidth}{0.4pt}
\renewcommand{\headrule}{\hbox to\headwidth{\color{RuleGray}\leaders\hrule height \headrulewidth\hfill}}

\newcommand{\dd}{\mathop{}\!\mathrm d}
\newcommand{\e}{\mathrm e}
\newcommand{\ii}{\mathrm i}
\newcommand{\vect}[1]{\bm{#1}}
\newcommand{\uvec}[1]{\hat{\vect{#1}}}
\newcommand{\dimens}[1]{\operatorname{dim}\!\left(#1\right)}
\newcommand{\order}[1]{\mathcal O\!\left(#1\right)}
\DeclareMathOperator{\proj}{proj}
\DeclareMathOperator{\sinc}{sinc}

\renewcommand{\theequation}{\thechapter.\arabic{equation}}
\renewcommand{\thefigure}{\thechapter.\arabic{figure}}
\renewcommand{\thetable}{\thechapter.\arabic{table}}
\makeatletter
\@addtoreset{equation}{chapter}
\@addtoreset{figure}{chapter}
\@addtoreset{table}{chapter}
\makeatother

\newcounter{bookexample}[chapter]
\renewcommand{\thebookexample}{\thechapter.\arabic{bookexample}}
\NewDocumentEnvironment{origexamplebox}{m}{%
  \refstepcounter{bookexample}%
  \begin{tcolorbox}[
    enhanced,breakable,colback=white,colframe=RuleGray,boxrule=0.5pt,arc=0mm,
    title={例~\thebookexample\quad #1},fonttitle=\bfseries,
    coltitle=Ink,colbacktitle=PaleBlue,borderline west={1.4pt}{0pt}{DeepBlue}]
}{\end{tcolorbox}}


% FinalReviewSolutionHeading20261024
\makeatletter
\newcommand{\LectureSolutionHeading}{%
  \par\noindent{\bfseries\color{DeepBlue}解答}\par\nobreak
  \@afterindentfalse\@afterheading}
\makeatother

\NewDocumentEnvironment{xesolution}{+b}{%
  \par\smallskip\begingroup\small\noindent
  \LectureSolutionHeading #1\par\endgroup\smallskip
}{}

\NewDocumentEnvironment{thinkanswer}{+b}{%
  \par\smallskip\begingroup\small\noindent
  {\bfseries\color{DeepBlue}思考辨析}\quad #1\par\endgroup\smallskip
}{}

% Used when a long worked solution is split by a page or a figure.  The
% continuation deliberately carries no second heading, so that it remains one
% answer both typographically and semantically.
\NewDocumentEnvironment{solutioncontinuation}{+b}{%
  \par\begingroup\small #1\par\endgroup\smallskip
}{}

% Numbering and restrained boxes for the independent advanced-reading volume.
% Advanced units deliberately do not consume the 32-chapter counter: their
% public identifiers are A1.1, A2.1, ... and remain stable whether a unit is
% compiled by itself or appended to the complete book.
\newcounter{advunit}
\newcounter{advexample}[advunit]
\newcounter{advdefinition}[advunit]
\newcounter{advtheorem}[advunit]
\renewcommand{\theadvexample}{A\arabic{advunit}.\arabic{advexample}}
\renewcommand{\theadvdefinition}{A\arabic{advunit}.\arabic{advdefinition}}
\renewcommand{\theadvtheorem}{A\arabic{advunit}.\arabic{advtheorem}}

\NewDocumentCommand{\AdvancedUnit}{m m}{%
  \clearpage
  \setcounter{advunit}{#1}%
  \setcounter{advexample}{0}%
  \setcounter{advdefinition}{0}%
  \setcounter{advtheorem}{0}%
  \setcounter{equation}{0}%
  \setcounter{figure}{0}%
  \setcounter{table}{0}%
  \setcounter{section}{0}%
  \setcounter{subsection}{0}%
  \renewcommand{\thesection}{A\arabic{advunit}.\arabic{section}}%
  \renewcommand{\thesubsection}{\thesection.\arabic{subsection}}%
  \renewcommand{\theequation}{A\arabic{advunit}.\arabic{equation}}%
  \renewcommand{\thefigure}{A\arabic{advunit}.\arabic{figure}}%
  \renewcommand{\thetable}{A\arabic{advunit}.\arabic{table}}%
  % Hyperref destinations must remain unique after the appendix and after
  % each advanced unit resets its visible section/figure counters.
  \renewcommand{\theHsection}{advA#1.section.\arabic{section}}%
  \renewcommand{\theHsubsection}{advA#1.section.\arabic{section}.\arabic{subsection}}%
  \renewcommand{\theHequation}{advA#1.equation.\arabic{equation}}%
  \renewcommand{\theHfigure}{advA#1.figure.\arabic{figure}}%
  \renewcommand{\theHtable}{advA#1.table.\arabic{table}}%
  \chapter*{\texorpdfstring{A#1\quad #2}{A#1 #2}}%
  \markboth{A#1\quad #2}{A#1\quad #2}%
  \addcontentsline{toc}{chapter}{\texorpdfstring{A#1\quad #2}{A#1 #2}}%
}

\NewDocumentEnvironment{advexamplebox}{m}{%
  \refstepcounter{advexample}%
  \begin{tcolorbox}[
    enhanced,breakable,colback=white,colframe=RuleGray,boxrule=0.5pt,arc=0mm,
    title={例~\theadvexample\quad #1},fonttitle=\bfseries,
    coltitle=Ink,colbacktitle=PaleBlue,borderline west={1.4pt}{0pt}{DeepBlue}]
}{\end{tcolorbox}}

\NewDocumentEnvironment{advsolution}{+b}{%
  \par\smallskip\begingroup\small\noindent
  \LectureSolutionHeading #1\par\endgroup\smallskip
}{}

\NewDocumentEnvironment{advdefinition}{m +b}{%
  \refstepcounter{advdefinition}%
  \begin{tcolorbox}[
    enhanced,breakable,colback=white,colframe=RuleGray,boxrule=0.5pt,arc=0mm,
    title={定义~\theadvdefinition\quad #1},fonttitle=\bfseries,
    coltitle=Ink,colbacktitle=PaleGold,borderline west={1.2pt}{0pt}{Gold}]
  #2
  \end{tcolorbox}
}{}

\NewDocumentEnvironment{advtheorem}{m +b}{%
  \refstepcounter{advtheorem}%
  \begin{tcolorbox}[
    enhanced,breakable,colback=white,colframe=RuleGray,boxrule=0.5pt,arc=0mm,
    title={定理~\theadvtheorem\quad #1},fonttitle=\bfseries,
    coltitle=Ink,colbacktitle=PaleBlue,borderline west={1.2pt}{0pt}{DeepBlue}]
  #2
  \end{tcolorbox}
}{}

\newtcolorbox{starread}[1]{
  enhanced,breakable,colback=PaleGold,colframe=RuleGray,boxrule=0.5pt,arc=0mm,
  title={选读：#1},fonttitle=\bfseries,coltitle=Ink,colbacktitle=PaleGold,
  borderline west={1.2pt}{0pt}{Gold}}

\newtcolorbox{methodbox}[1]{
  enhanced,breakable,colback=PaleBlue,colframe=RuleGray,boxrule=0.5pt,arc=0mm,
  title={方法：#1},fonttitle=\bfseries,coltitle=Ink,colbacktitle=PaleBlue,
  borderline west={1.2pt}{0pt}{DeepBlue}}
\newtcolorbox{warningbox}[1]{
  enhanced,breakable,colback=SoftGray,colframe=RuleGray,boxrule=0.5pt,arc=0mm,
  title={易错提醒：#1},fonttitle=\bfseries,coltitle=Ink,colbacktitle=SoftGray,
  borderline west={1.2pt}{0pt}{TextGray}}
\newtcolorbox{extensionbox}[1]{
  enhanced,breakable,colback=PaleGold,colframe=RuleGray,boxrule=0.5pt,arc=0mm,
  title={拓展：#1},fonttitle=\bfseries,coltitle=Ink,colbacktitle=PaleGold,
  borderline west={1.2pt}{0pt}{Gold}}

\newcommand{\Needspace}[1]{%
  \par\ifdim\dimexpr\pagegoal-\pagetotal\relax<#1\relax\newpage\fi}

\newcommand{\bookcaption}[1]{%
  \refstepcounter{figure}%
  \par\nopagebreak[4]\smallskip
  {\centering\small\color{TextGray}图~\thefigure：#1\par}\smallskip}
\newcommand{\booktablecaption}[1]{%
  \refstepcounter{table}%
  \par\nopagebreak[4]\smallskip
  {\centering\small\color{TextGray}表~\thetable：#1\par}\smallskip}
\newcommand{\diagramcaption}[1]{\bookcaption{#1}}
\newcommand{\figcap}[1]{\bookcaption{#1}}
\newcommand{\origcaption}[1]{\bookcaption{#1}}
% Keep newly authored diagrams and their complete captions on the same page.
\NewDocumentEnvironment{qualityfigure}{}{%
  \par\begin{center}\begin{minipage}{\linewidth}\centering
}{\end{minipage}\end{center}}
\newcommand{\historyfigure}[3][0.9\linewidth]{%
  \begin{center}\begin{minipage}{#1}\centering
  \includegraphics[width=\linewidth]{#2}\par\smallskip
  \bookcaption{#3}
  \end{minipage}\end{center}}
\newcommand{\sourceimage}[3][0.9\linewidth]{\historyfigure[#1]{#2}{#3}}
\newcommand{\redrawimage}[2][0.72\linewidth]{%
  \begin{center}\includegraphics[width=#1]{#2}\end{center}}
\newcommand{\unitheading}[2]{}
\newcommand{\examplesource}[1]{}

\tikzset{
  every node/.append style={fill=none,inner sep=1.2pt},
  clearlabel/.style={fill=none,inner sep=1pt,text=Ink},
  point/.style={circle,fill=Ink,draw=none,inner sep=1.7pt},
  vector/.style={-{Latex[length=2.5mm]},very thick,draw=DeepBlue},
  component/.style={-{Latex[length=2.2mm]},thick,draw=TextGray},
  axis/.style={-{Latex[length=2mm]},thin,draw=TextGray},
  guide/.style={densely dashed,draw=TextGray},
  curve/.style={very thick,draw=DeepBlue},
  ground/.style={very thick,draw=Ink},
  spring/.style={decorate,decoration={zigzag,segment length=4pt,amplitude=2pt},thick,draw=DeepBlue},
  tangent/.style={thick,draw=TextGray},
  secant/.style={thick,draw=DeepBlue},
  kaxis/.style={-{Latex[length=2mm]},thin,draw=TextGray},
  kcurve/.style={very thick,draw=DeepBlue},
  kguide/.style={densely dashed,draw=TextGray}
}
